% ============================================================
% Chapter 5: Development
%   Section: News Acquisition: Discovery and Extraction
%     Subsection: Reading Evidence Pages
% Included from chapters/05-development/02-news-acquisition/news-acquisition.tex
% ============================================================

\subsection{Reading Evidence Pages}
\label{subsec:scraping-evidence}

The same extractor reads the web pages that evidence retrieval finds
for each claim (\path{services/fact_checker/retrieval/scraper.py}),
rather than a second scraping stack. A search result has no configured
source, so it is extracted with a generic placeholder source (neutral
reliability, never stored), and the extractor still matches its domain
against the catalogue, so a hit on a configured outlet is read with
that outlet's selectors. Three things differ from article extraction:

\begin{itemize}
    \item \textbf{No browser and no publication time.} Evidence is the
    most frequent purpose and already has a fallback, so the expensive
    steps are left out (Section~\ref{subsec:scraping-cascade}).
    \item \textbf{A failed page keeps its search result.} When a page
    cannot be read, the evidence keeps the title and snippet the search
    engine returned, and the claim is judged on that rather than losing
    the source.
    \item \textbf{Concurrent, and in input order.} A claim's pages are
    fetched at once instead of one after another; this was the
    pipeline's longest step apart from the language model, with one
    slow host holding up the others for no reason. A process-wide
    ceiling of eight concurrent fetches (\texttt{SCRAPE\_MAX\_CONCURRENCY})
    is held around each fetch only, because several claims and several
    articles reach this point at once. The pages come back in the
    order they were requested, not the order they finished: the list
    is ranked, cut and handed to the language model as a numbered list
    that it cites by index, so completion order would silently point
    every citation at a different source
    (Section~\ref{sec:design-concurrency}).
\end{itemize}

How many search results are read per claim is itself a per-run
threshold (\path{evidence_fetch_candidates}, eight by default),
because widening recall is the first thing to try when a claim comes
back \texttt{UNVERIFIED} for lack of evidence.
